\documentclass[reqno]{amsart}
\usepackage[T1]{fontenc}
\usepackage{amsmath,amssymb,amsthm}
\usepackage{xr-hyper}
\usepackage[hidelinks]{hyperref}
\externaldocument[da-]{20261126T030300--the-data__workbook}
\externaldocument[ob-]{20261126T040100--spaces-forms-and-twists__workbook}
\externaldocument[tn-]{20261125T002100--definite-theta-lifts-and-norms__workbook}
\externaldocument[fm-]{20261125T009500--fock-model-and-theta-intertwiners__workbook}
\externaldocument[ss-]{20261126T030400--the-seesaw-identity__workbook}
\externaldocument[s1-]{20261126T051200--transferring-a-bessel-period-through-a-theta-lift__workbook}
\externaldocument[s3-]{20261126T051400--lattices-cuts-and-the-ramified-forms__workbook}

\newcounter{chap}
% Chapter numbers, generated from 20261125T012000--grh-workbook__guide.org
% by scripts/make-chapters.py.  Do not edit: change the order of the
% index and rerun the script (./preview-site.sh and CI do this).
%
% Usage in a sheet, after \newcounter{chap}:
%   % Chapter numbers, generated from 20261125T012000--grh-workbook__guide.org
% by scripts/make-chapters.py.  Do not edit: change the order of the
% index and rerun the script (./preview-site.sh and CI do this).
%
% Usage in a sheet, after \newcounter{chap}:
%   % Chapter numbers, generated from 20261125T012000--grh-workbook__guide.org
% by scripts/make-chapters.py.  Do not edit: change the order of the
% index and rerun the script (./preview-site.sh and CI do this).
%
% Usage in a sheet, after \newcounter{chap}:
%   \input{chapters}\chapnumber{<sheet ID>}
\makeatletter
\providecommand\chapters@set[2]{\@namedef{chapters@#1}{#2}}
\chapters@set{20261125T001000}{1}% zeros-to-ramp-sums
\chapters@set{20261125T001100}{2}% inert-primes-and-estimate-b
\chapters@set{20261125T005000}{3}% galois-fourier-and-averaging
\chapters@set{20261125T005500}{4}% columns-from-a-sphere
\chapters@set{20261125T000000}{5}% slope-method
\chapters@set{20261126T030100}{6}% a-twist-for-each-prime
\chapters@set{20261125T003000}{7}% the-chain-and-the-square-root
\chapters@set{20261125T010000}{8}% rank-lemma-and-zero-estimate
\chapters@set{20261125T010500}{9}% poincare-series-and-near-orthogonality
\chapters@set{20261124T223000}{10}% elliptic-curves-for-the-proof
\chapters@set{20261125T009500}{11}% fock-model-and-theta-intertwiners
\chapters@set{20261124T223100}{12}% complex-multiplication-and-inert-reduction
\chapters@set{20261125T008000}{13}% one-cm-determinant
\chapters@set{20261124T214500}{14}% supersingular-theta-divisibility
\chapters@set{20261124T214600}{15}% the-isotropic-plane
\chapters@set{20261125T013700}{16}% pairings-on-a-formal-group-of-height-two
\chapters@set{20261125T013500}{17}% hecke-nondegeneracy-on-the-tree
\chapters@set{20261126T030200}{18}% the-bound-is-attained
\chapters@set{20261125T009000}{19}% kronecker-limit-and-chowla-selberg
\chapters@set{20261125T002000}{20}% definite-theta-lifts
\chapters@set{20261125T015500}{21}% siegel-weil-in-the-smallest-case
\chapters@set{20261125T002100}{22}% definite-theta-lifts-and-norms
\chapters@set{20261125T016500}{23}% hecke-operators-and-theta-lifts
\chapters@set{20261125T016700}{24}% divisibility-along-the-chain
\chapters@set{20261125T004000}{25}% denominators-and-sha
\chapters@set{20261125T004100}{26}% tate-twists-and-sha
\chapters@set{20261126T020430}{27}% the-minus-part-by-gauss-sums
\chapters@set{20261125T014100}{28}% congruence-modules
\chapters@set{20261125T014500}{29}% eisenstein-congruences-and-selmer-classes
\chapters@set{20261125T014700}{30}% pseudocharacters-and-the-length-bound
\chapters@set{20261125T011000}{31}% galois-representations-black-box
\chapters@set{20261126T030300}{32}% the-data
\chapters@set{20261126T020210}{33}% bessel-poincare-series-and-their-gram-operators
\chapters@set{20261126T030400}{34}% the-seesaw-identity
\chapters@set{20261126T030500}{35}% divisibility-of-the-composite-lift
\chapters@set{20261126T030600}{36}% theorem-a-short-forms
\chapters@set{20261126T030700}{37}% the-determinant-of-a-chain
\chapters@set{20261126T030800}{38}% the-cost-of-one-lift
\chapters@set{20261126T030900}{39}% denominators-of-projections-to-lifts
\chapters@set{20261126T031000}{40}% twists-with-few-congruences
\chapters@set{20261126T031100}{41}% theorem-c-and-the-main-inequality
\providecommand\chapnumber[1]{%
  \@ifundefined{chapters@#1}%
    {\PackageError{chapters}{Sheet #1 is not listed in the index}%
      {List it in 20261125T012000--grh-workbook__guide.org\MessageBreak
       and run scripts/make-chapters.py.}}%
    {\edef\thechap{\@nameuse{chapters@#1}}}}
\makeatother
\chapnumber{<sheet ID>}
\makeatletter
\providecommand\chapters@set[2]{\@namedef{chapters@#1}{#2}}
\chapters@set{20261125T001000}{1}% zeros-to-ramp-sums
\chapters@set{20261125T001100}{2}% inert-primes-and-estimate-b
\chapters@set{20261125T005000}{3}% galois-fourier-and-averaging
\chapters@set{20261125T005500}{4}% columns-from-a-sphere
\chapters@set{20261125T000000}{5}% slope-method
\chapters@set{20261126T030100}{6}% a-twist-for-each-prime
\chapters@set{20261125T003000}{7}% the-chain-and-the-square-root
\chapters@set{20261125T010000}{8}% rank-lemma-and-zero-estimate
\chapters@set{20261125T010500}{9}% poincare-series-and-near-orthogonality
\chapters@set{20261124T223000}{10}% elliptic-curves-for-the-proof
\chapters@set{20261125T009500}{11}% fock-model-and-theta-intertwiners
\chapters@set{20261124T223100}{12}% complex-multiplication-and-inert-reduction
\chapters@set{20261125T008000}{13}% one-cm-determinant
\chapters@set{20261124T214500}{14}% supersingular-theta-divisibility
\chapters@set{20261124T214600}{15}% the-isotropic-plane
\chapters@set{20261125T013700}{16}% pairings-on-a-formal-group-of-height-two
\chapters@set{20261125T013500}{17}% hecke-nondegeneracy-on-the-tree
\chapters@set{20261126T030200}{18}% the-bound-is-attained
\chapters@set{20261125T009000}{19}% kronecker-limit-and-chowla-selberg
\chapters@set{20261125T002000}{20}% definite-theta-lifts
\chapters@set{20261125T015500}{21}% siegel-weil-in-the-smallest-case
\chapters@set{20261125T002100}{22}% definite-theta-lifts-and-norms
\chapters@set{20261125T016500}{23}% hecke-operators-and-theta-lifts
\chapters@set{20261125T016700}{24}% divisibility-along-the-chain
\chapters@set{20261125T004000}{25}% denominators-and-sha
\chapters@set{20261125T004100}{26}% tate-twists-and-sha
\chapters@set{20261126T020430}{27}% the-minus-part-by-gauss-sums
\chapters@set{20261125T014100}{28}% congruence-modules
\chapters@set{20261125T014500}{29}% eisenstein-congruences-and-selmer-classes
\chapters@set{20261125T014700}{30}% pseudocharacters-and-the-length-bound
\chapters@set{20261125T011000}{31}% galois-representations-black-box
\chapters@set{20261126T030300}{32}% the-data
\chapters@set{20261126T020210}{33}% bessel-poincare-series-and-their-gram-operators
\chapters@set{20261126T030400}{34}% the-seesaw-identity
\chapters@set{20261126T030500}{35}% divisibility-of-the-composite-lift
\chapters@set{20261126T030600}{36}% theorem-a-short-forms
\chapters@set{20261126T030700}{37}% the-determinant-of-a-chain
\chapters@set{20261126T030800}{38}% the-cost-of-one-lift
\chapters@set{20261126T030900}{39}% denominators-of-projections-to-lifts
\chapters@set{20261126T031000}{40}% twists-with-few-congruences
\chapters@set{20261126T031100}{41}% theorem-c-and-the-main-inequality
\providecommand\chapnumber[1]{%
  \@ifundefined{chapters@#1}%
    {\PackageError{chapters}{Sheet #1 is not listed in the index}%
      {List it in 20261125T012000--grh-workbook__guide.org\MessageBreak
       and run scripts/make-chapters.py.}}%
    {\edef\thechap{\@nameuse{chapters@#1}}}}
\makeatother
\chapnumber{<sheet ID>}
\makeatletter
\providecommand\chapters@set[2]{\@namedef{chapters@#1}{#2}}
\chapters@set{20261125T001000}{1}% zeros-to-ramp-sums
\chapters@set{20261125T001100}{2}% inert-primes-and-estimate-b
\chapters@set{20261125T005000}{3}% galois-fourier-and-averaging
\chapters@set{20261125T005500}{4}% columns-from-a-sphere
\chapters@set{20261125T000000}{5}% slope-method
\chapters@set{20261126T030100}{6}% a-twist-for-each-prime
\chapters@set{20261125T003000}{7}% the-chain-and-the-square-root
\chapters@set{20261125T010000}{8}% rank-lemma-and-zero-estimate
\chapters@set{20261125T010500}{9}% poincare-series-and-near-orthogonality
\chapters@set{20261124T223000}{10}% elliptic-curves-for-the-proof
\chapters@set{20261125T009500}{11}% fock-model-and-theta-intertwiners
\chapters@set{20261124T223100}{12}% complex-multiplication-and-inert-reduction
\chapters@set{20261125T008000}{13}% one-cm-determinant
\chapters@set{20261124T214500}{14}% supersingular-theta-divisibility
\chapters@set{20261124T214600}{15}% the-isotropic-plane
\chapters@set{20261125T013700}{16}% pairings-on-a-formal-group-of-height-two
\chapters@set{20261125T013500}{17}% hecke-nondegeneracy-on-the-tree
\chapters@set{20261126T030200}{18}% the-bound-is-attained
\chapters@set{20261125T009000}{19}% kronecker-limit-and-chowla-selberg
\chapters@set{20261125T002000}{20}% definite-theta-lifts
\chapters@set{20261125T015500}{21}% siegel-weil-in-the-smallest-case
\chapters@set{20261125T002100}{22}% definite-theta-lifts-and-norms
\chapters@set{20261125T016500}{23}% hecke-operators-and-theta-lifts
\chapters@set{20261125T016700}{24}% divisibility-along-the-chain
\chapters@set{20261125T004000}{25}% denominators-and-sha
\chapters@set{20261125T004100}{26}% tate-twists-and-sha
\chapters@set{20261126T020430}{27}% the-minus-part-by-gauss-sums
\chapters@set{20261125T014100}{28}% congruence-modules
\chapters@set{20261125T014500}{29}% eisenstein-congruences-and-selmer-classes
\chapters@set{20261125T014700}{30}% pseudocharacters-and-the-length-bound
\chapters@set{20261125T011000}{31}% galois-representations-black-box
\chapters@set{20261126T030300}{32}% the-data
\chapters@set{20261126T020210}{33}% bessel-poincare-series-and-their-gram-operators
\chapters@set{20261126T030400}{34}% the-seesaw-identity
\chapters@set{20261126T030500}{35}% divisibility-of-the-composite-lift
\chapters@set{20261126T030600}{36}% theorem-a-short-forms
\chapters@set{20261126T030700}{37}% the-determinant-of-a-chain
\chapters@set{20261126T030800}{38}% the-cost-of-one-lift
\chapters@set{20261126T030900}{39}% denominators-of-projections-to-lifts
\chapters@set{20261126T031000}{40}% twists-with-few-congruences
\chapters@set{20261126T031100}{41}% theorem-c-and-the-main-inequality
\providecommand\chapnumber[1]{%
  \@ifundefined{chapters@#1}%
    {\PackageError{chapters}{Sheet #1 is not listed in the index}%
      {List it in 20261125T012000--grh-workbook__guide.org\MessageBreak
       and run scripts/make-chapters.py.}}%
    {\edef\thechap{\@nameuse{chapters@#1}}}}
\makeatother
\chapnumber{20261126T051300}
\renewcommand{\thesection}{\thechap.\arabic{section}}

\theoremstyle{definition}
\newtheorem{problem}{Problem}[section]
\newtheorem{definition}[problem]{Definition}
\theoremstyle{plain}
\newtheorem{theorem}{Theorem}[section]
\theoremstyle{remark}
\newtheorem*{leanremark}{In the Lean formalization}

\newcommand{\A}{\mathbb{A}}
\newcommand{\ii}{\mathrm{i}}
\newcommand{\Ac}{\mathcal{A}}
\newcommand{\Uc}{\mathcal{U}}
\newcommand{\Wc}{\mathcal{W}}
\newcommand{\Hom}{\operatorname{Hom}}
\newcommand{\Heis}{\mathrm{Heis}}
\newcommand{\st}{\mathrm{st}}
\newcommand{\CT}{\mathrm{CT}}
\newcommand{\zt}{\mathrm{zt}}
\newcommand{\ev}{\mathrm{ev}}
\newcommand{\coeff}{\operatorname{coeff}}
\newcommand{\diag}{\operatorname{diag}}
\newcommand{\len}{\operatorname{len}}
\newcommand{\fm}{f_{\mathrm{mid}}}
\newcommand{\Wtr}{\Wc^{\mathrm{tr}}}
\newcommand{\bu}{\mathbf u}
\newcommand{\tl}[1]{\texttt{#1}}

\setlength{\emergencystretch}{3em}

\title{Two lifts at a time}
\author{}
\date{November 26, 2026}

\begin{document}
\maketitle
\refstepcounter{chap}\label{chap}

\paragraph{Prerequisites.}
Chapter~\ref{s1-chap}: spaces with a string, grades and degrees, string maps, constant terms, the transfer (Definitions~\ref{s1-def:string}--\ref{s1-def:transfer}), the Bessel groups \(R = \Uc \rtimes U(X)\), \(\psi_E\), \(j_E\), \(\psi_\Uc\), and Theorem~\ref{s1-thm:transfer}, in its notation.  Chapter~\ref{ss-chap}, Definitions~\ref{ss-def:lower} and~\ref{ss-def:definite}: the Bessel function \(\Wc_{\lambda,\varphi}\); lift \(j\) of the lower chain \(\Theta_{\xi,j}\), its Weil representation \(\omega_j\), its datum \(\phi_j(\xi)\) (a harmonic at infinity times the finite datum \(\mathbf 1_j(\xi)\)), the levels \(K^Y_j(\xi) \subseteq K^{Y\circ}_j\); Heisenberg translations \(\Heis\); the definite lifts \(\Theta^X_{\xi,j}(v)\).  Chapter~\ref{tn-chap}: Problems~\ref{tn-pr:U1U1}, \ref{tn-pr:factor}, \ref{tn-pr:compose}.  Chapters~\ref{da-chap}, \ref{ob-chap}: \(X_j\), \(Y_j\), \(\varrho = m/2 - 1\), tame characters (of weight zero), Problem~\ref{da-pr:piN}.

In the proof of the seesaw identity (Chapter~\ref{ss-chap}) the Bessel period of \(\varphi\) is transferred along the lower chain one lift at a time (Theorem~\ref{s1-thm:transfer}).  How does the function left on \(Y_j\) become a lift between the \emph{definite} spaces \(X_j\)?  The proof groups the transfers in pairs: the string maps of two consecutive lifts fit into one isotropic map.  We compute it for strings of lengths \(3, 2, 1\); then, after each pair, we identify the translated stage function on the definite group.

Lift \(j\) (\(1 \le j \le m\)) goes from \(Y_{j-1}\) to \(Y_j\); pair \(i\) (\(0 \le i \le \varrho\)) consists of lifts \(2i+1\), \(2i+2\).  For \(j < m\), \(Y_j\) is \(X_j\) plus a string of length \(m-1-j\), with Bessel group \(R_j\); \(Y_m\) is definite, of rank \(n_m - 1\), with string of length \(0\) and \(R_m = U(Y_m)\) (Chapter~\ref{da-chap}).


\section{Stage functions}\label{sec:stage}

Fix \(\lambda\), archimedean points \(g_j\) of \(U(Y_j)\), finite points \(y_j \in U(Y_j)(\A_f)\) (\(j \ge 0\)), a tame character \(\xi\) and \(\varphi \in S\).

\begin{definition}\label{def:stage}
Put \(\Wtr_0 := \Wc_{\lambda,\varphi}\) on \(R_0(\A)\).  For \(0 \le j \le m - 1\) let
\[
  \Wc_{j+1}(h) := \int_{[R_j]} \Wtr_j(r)\,\CT_{\st_{j+1}}\bigl(\omega_{j+1}(r\,(g_j, y_j), h)\,\phi_{j+1}(\xi)\bigr)\,dr
\]
on \(U(Y_{j+1})(\A)\), the transfer of \(\Wtr_j\) for \(\omega_{j+1}((g_j, y_j), 1)\,\phi_{j+1}(\xi)\), and, for \(j + 1 < m\), \(\Wtr_{j+1}(r) := \Wc_{j+1}(r\,(g_{j+1}, y_{j+1}))\) on \(R_{j+1}(\A)\).  These are the \emph{stage functions}.
\end{definition}

In the proof, for \(0 \le j < m\), \(\Wtr_j\) satisfies the hypothesis of Definition~\ref{s1-def:transfer} (continuous, left \(R_j(F)\)-invariant, \(\Wtr_j(nr) = \psi_\Uc(n)\Wtr_j(r)\) when \(n - 1\) has degree \(\ge 2\)); call such functions \emph{equivariant}.  A \emph{form of the conjugate type} on \(Y_j\) is defined as a form of the lower chain on \(Y_j\), with the type kernel replaced by its complex conjugate; these forms are paired with those of the lower chain.  In the proof, for \(0 \le j < m\) there is a linear map \(\Theta^{\top}_{\xi,j+1}\), the \emph{transpose} of lift \(j+1\), from the forms of the conjugate type on \(Y_{j+1}\) to those on \(Y_j\), with
\[
  \int_{[U(Y_{j+1})]} z\;\Theta_{\xi,j+1}f = \int_{[U(Y_j)]} (\Theta^{\top}_{\xi,j+1}z)\,f
\]
for every form \(f\) of the lower chain on \(Y_j\) and every form \(z\) of the conjugate type on \(Y_{j+1}\); and for such \(z\),
\[
  (\Theta_{\xi,j+1}^{\top} z)(g) = [K^{Y\circ}_j : K^Y_j(\xi)] \int_{[U(Y_{j+1})]} \theta\bigl(\omega_{j+1}(g, h)\,\phi_{j+1}(\xi)\bigr)\, z(h)\,dh ,
\]
and \(z\) is automorphic and rapidly decreasing.

\begin{problem}\label{pr:stage}
For \(0 \le j \le m-1\) and every form \(z\) of the conjugate type on \(Y_{j+1}\), the integral on the right converges absolutely and
\[
  \int_{[R_j]} \Wtr_j(r)\,(\Theta^{\top}_{\xi,j+1} z)\bigl(r\,(g_j, y_j)\bigr)\,dr = [K^{Y\circ}_j : K^Y_j(\xi)] \int_{R_{j+1}(F)\backslash U(Y_{j+1})(\A)} \Wc_{j+1}\, z .
\]
\end{problem}

Chapter~\ref{s3-chap} uses this identity in the recursion for paired periods.


\section{Three strings}\label{sec:strings}

For adelic maps and data we use the CM field, Weil data and framed spaces of Chapter~\ref{s1-chap}.  Write \(V^-\) for \(V\) with its form negated.

\begin{definition}\label{def:stacked}
Let \(V, V', V''\) be spaces with strings for the same \(N\), of lengths \(\ell+1, \ell, \max(\ell-1, 0)\) (\(\ell \ge 0\)), with string maps \(\st_1 \colon V \to V'\) and \(\st_2 \colon V' \to V''\).  The \emph{stacked string map} is \(A := \binom{\st_1^*}{\st_2} \colon V' \to V^- \oplus V''\).
\end{definition}

In Problems~\ref{pr:st-instance}--\ref{pr:adjoint-degree} take \(E = \mathbb Q(\ii)\), \(N = 5\), \(c_V = c_{V'} = c_{V''} = 1\), no definite parts, \(J_V = \mathrm{antidiag}(1, -5, 1)\), \(J_{V'} = \mathrm{antidiag}(1, 1)\), \(J_{V''} = (-5)\).

\begin{problem}\label{pr:st-instance}
\(\st_1 = \bigl(\begin{smallmatrix} 1&0&0\\ 0&-5&0 \end{smallmatrix}\bigr)\), \(\st_2 = (1, 0)\), and \(\st_1\st_1^* = \st_2^*\st_2 = \bigl(\begin{smallmatrix} 0&0\\ -5&0 \end{smallmatrix}\bigr)\).
\end{problem}

\begin{problem}\label{pr:A-instance}
\(A\) has rows \((0,0)\), \((1,0)\), \((0,1)\), \((1,0)\), and \(A^*A = 0\) for the form \(-J_V \oplus J_{V''}\).
\end{problem}

\begin{problem}\label{pr:adjoint-degree}
For the grades of \(V'\) and of \(V\) as weights, the maps \(a \colon V \to V'\) of positive degree are the \(a = \bigl(\begin{smallmatrix} 0&\alpha_0&\alpha_1\\ 0&0&\alpha_2 \end{smallmatrix}\bigr)\); then \(a^* = \bigl(\begin{smallmatrix} \bar\alpha_2&\bar\alpha_1\\ 0&-\bar\alpha_0/5\\ 0&0 \end{smallmatrix}\bigr)\), and \(a \mapsto a^*\) is a conjugate-linear bijection onto the maps \(V' \to V\) of positive degree.
\end{problem}

\begin{problem}\label{pr:even-instance}
Over \(\mathbb Q(\ii)\), for three spaces with no definite parts, string lengths \(4, 3, 2\), \(N = 5\) and \(c_V = c_{V'} = c_{V''} = 2\): \(\st_1 = \bigl(\begin{smallmatrix} 2&0&0&0\\ 0&2&0&0\\ 0&0&1&0 \end{smallmatrix}\bigr)\), \(\st_2 = \bigl(\begin{smallmatrix} 1&0&0\\ 0&-10&0 \end{smallmatrix}\bigr)\), and \(\st_1\st_1^* = \st_2^*\st_2 = \bigl(\begin{smallmatrix} 0&0&0\\ 2&0&0\\ 0&-10&0 \end{smallmatrix}\bigr)\).
\end{problem}

\begin{problem}\label{pr:isotropic}
Let \(\ell \ge 1\) and \(c_V = c_{V'}\).  Then \(\st_1\st_1^* = \st_2^*\st_2\); this map kills the definite part and \(f'_{\ell-1}\) and sends \(f'_q\) to a multiple of \(f'_{q+1}\) for \(q < \ell - 1\).  Also \(A^*A = -\st_1\st_1^* + \st_2^*\st_2 = 0\).
\end{problem}

In pair \(i\) the strings of \(Y_{2i}, Y_{2i+1}, Y_{2i+2}\) have lengths \(2(\varrho - i) + 1\), \(2(\varrho - i)\), \(2(\varrho - i) - 1\) if \(i < \varrho\) (a \emph{middle pair}) and \(1, 0, 0\) if \(i = \varrho\) (the \emph{last pair}).  The gradings of \(Y_j\), \(j < m\), have one common constant; that of \(Y_m\) has constant \(1\).  In a middle pair the first string has odd length, so a map \(V' \to V^- \oplus V''\) homogeneous of degree zero kills the string of \(V'\) and maps its definite part into the span of the definite parts and the middle vectors \(\fm\) of \(V\) and \(\fm'\) of \(V''\), of norms \(N\) in \(V^-\) and \(-N\) in \(V''\).

\begin{problem}\label{pr:plane}
For norms \(5\), \(-5\), the basis \(\tfrac1{10}(\fm - \fm')\), \(\fm + \fm'\) of the plane they span has Gram matrix \(\bigl(\begin{smallmatrix} 0&1\\ 1&0 \end{smallmatrix}\bigr)\); in general so has \((\fm - \fm')/2N\), \(\fm + \fm'\).
\end{problem}

So in a middle pair the part of grade zero of \(Y_{2i}^- \oplus Y_{2i+2}\) is \(X_{2i}^- \oplus X_{2i+2}\) plus one hyperbolic plane.

\begin{problem}\label{pr:last-pair}
In the last pair, with \(\fm\) the vector of the string of \(Y_{2\varrho}\), \(\fm/\pi_N\) has norm \(1\) in \(Y_{2\varrho}^-\), for \(N = 5\), \(\pi_N = 2 + \ii\) and in general.  The part of grade zero of \(Y_{2\varrho}^- \oplus Y_m\) is \(X_{2\varrho}^- \oplus X_m\), no plane being left, with \(X_m = Y_m \perp E\,\fm/\pi_N\), and \(h \in U(Y_m)\) acts on \(X_m\) as \(\diag(h, 1)\).
\end{problem}

\begin{problem}\label{pr:A-last}
In the last pair \(A = 0\).
\end{problem}

In the proof, \(\Heis(w)\Heis(w') = \psi_E\bigl(\tfrac12 j_E \operatorname{tr}(w'^*w)\bigr)\Heis(w + w')\) for adelic maps \(w, w'\) between two spaces, the adjoint taken for their forms; we use it for \(\Hom(V', V^- \oplus V'')\).

\begin{problem}\label{pr:commute}
Under the hypotheses of Problem~\ref{pr:isotropic}, \((An')^*(An) = 0\) and \(\Heis(An)\Heis(An') = \Heis(A(n + n'))\) for all adelic endomorphisms \(n, n'\) of \(V'\).
\end{problem}


\section{Two constant terms at once}\label{sec:two-ct}

Let \(V, V', V''\) be spaces with strings for the same \(N\), of lengths \(\ell+1, \ell, \max(\ell-1, 0)\), \(\ell \ge 0\) even, \(c_V = c_{V'}\).  For the splitting into definite spaces and the two-period identity take \(V = Y_{2i}\), \(V' = Y_{2i+1}\), \(V'' = Y_{2i+2}\) with the middle or last double-lift datum of pair \(i\).  Write \(X, X', X''_0\) for the definite parts and \(R, R'\) for the Bessel groups of \(V, V'\).  For data \(\phi\) on \(\Hom(V, V')\) and \(\phi'\) on \(\Hom(V', V'')\), the \emph{join} \(\phi \boxtimes \phi'\) on \(\Hom(V', V^- \oplus V'')\) is, in the proof, the product of \(\phi'\) with \(\phi\) carried to \(\Hom(V', V^-)\) by a scaled adjoint change of coordinates; for all adelic \(w, w'\), with matrices in the rational bases, \(\Heis\binom{w^*}{w'}(\phi \boxtimes \phi') = \Heis(w)\phi \boxtimes \Heis(w')\phi'\).  We use:

(i) \emph{Product formula}, in the proof: for framed spaces \(V, V', V''\) (definite ones included) there is \(z_\theta \ne 0\) with \(\theta(\phi \boxtimes \phi') = z_\theta\,\theta(\phi)\theta(\phi')\) for all \(\phi, \phi'\).  Compare Problem~\ref{tn-pr:factor}; for the Weil representation see \cite{We}.

For data on \(\Hom(V', V^- \oplus V'')\) let \(\zt_{>0}(\cdot)(0)\) be the mean of \(\theta(\Heis(b)\,\cdot\,)\) over the adelic maps \(b \colon V' \to V^- \oplus V''\) of positive degree modulo the rational ones (the weights on \(V^- \oplus V''\) are the grades); at a point \(v\) of the retained grade-zero space, \(\zt_{>0}(\phi)(v)\) is the same mean of \(\theta(\Heis(\iota v + b)\phi)\), \(\iota\) the inclusion of that space.

\begin{problem}\label{pr:ct-product}
For adelic maps \(w \colon V \to V'\) and \(w' \colon V' \to V''\),
\[
  \CT_w \phi \cdot \CT_{w'} \phi' = z_\theta^{-1}\,\zt_{>0}\bigl(\Heis(w)\phi \boxtimes \Heis(w')\phi'\bigr)(0) .
\]
\end{problem}

\begin{problem}\label{pr:odd}
If \(\len V\) is odd, all grades of \(V\) are even, and \(n - 1\) has degree \(\ge 2\) for every \(n \in \Uc(\A)\).
\end{problem}

In the proof, for the first edge with the Weil action \(\omega\) of the rank twists of an automorphic character of the norm-one torus, \(\Wc(nr)\,\CT_{\st_1}(\omega(nr, h)\phi) = \Wc(r)\,\CT_{\st_1}(\omega(r, h)\phi)\) for every equivariant \(\Wc\), \(n \in \Uc(\A)\), \(r \in R(\A)\), \(h \in U(V')(\A)\).  Let \(\zt\) be the zero term along the maps of positive degree followed, in a middle pair, by the mean over the adelic maps modulo the rational maps from \(X'\) into the line \(E(\fm + \fm')\) of Problem~\ref{pr:plane} (no plane in the last pair).  It lands in data on \(\Hom(X', X^- \oplus X'')\), where \(X'' = X''_0\) in a middle pair and \(X'' = X_m\) in the last pair.

(iv) \emph{Bessel groups}, in the proof.  Let \(R = \Uc \rtimes U(X)\) be the Bessel group of a space \(W\) with a string; \([R]\) is compact.  Let \(b_1, \dots, b_e\) be distinct basis vectors of grade one (\(e \ge 0\); none if \(\len W\) is odd) and \(n^{(1)} := (c_W (n - 1)_{b_p, x})_{p,x}\), \(x\) running over the basis of \(X\).  There is \(c_R\) such that: if \(H\) is continuous and left \(R(F)\)-invariant on \(R(\A)\), \(H(n\gamma) = G_\gamma(n^{(1)})\) for \(n \in \Uc(\A)\), \(\gamma \in U(X)(\A)\), each \(G_\gamma\) is continuous and periodic under rational matrices, and its mean \(\langle G_\gamma\rangle\) is continuous and left \(U(X)(F)\)-invariant, then \(\int_{[R]} H = c_R \int_{[U(X)]} \langle G_\gamma\rangle\,d\gamma\).  See \cite{Bo}.

For a middle or last pair we use the following two-period identity from the proof.  Let \(\omega\), \(\omega'\) be the Weil representations of the two edges, with automorphic characters \(\mu\), \(\mu'\) of the norm-one torus, each multiplied on a space of rank \(d\) by \(\upsilon^{\lfloor d/2\rfloor}\).  For data \(\phi, \phi'\), an equivariant \(\Wc\) on \(R(\A)\) and \(h \in U(X''_0)(\A)\), acting on \(V''\) through \(X''_0\),
\begin{multline}\label{eq:two-periods}
  \int_{[R']}\!\int_{[R]} \Wc(r)\,\CT_{\st_1}\bigl(\omega(r, r')\phi\bigr)\,\CT_{\st_2}\bigl(\omega'(r', h)\phi'\bigr)\,dr\,dr' \\
  = \frac{c_R c_{R'}}{z_\theta} \int_{[U(X)]} \Wc(\gamma) \int_{[U(X')]} \theta\bigl(\omega^\triangle(\gamma, \gamma', \bar h)\,\zt\,\Heis(A)(\phi \boxtimes \phi')\bigr)\,d\gamma'\,d\gamma ,
\end{multline}
where \(c_R, c_{R'}\) do not depend on \(\mu, \mu', \phi, \phi', \Wc\) or \(h\); \(\omega^\triangle\) is the Weil representation of \(U(X) \times U(X') \times U(X'')\) on \(\Hom(X', X^- \oplus X'')\) normalized as in the proof (the characters \(\mu\), \(\mu'\) and a flip character); and \(\bar h = h\) in a middle pair, \(\bar h = \diag(h, 1)\) in the last.


\section{The definite double lift}\label{sec:definite}

The right side of \eqref{eq:two-periods} is a double theta integral between definite spaces.  First the Gaussian case.

\begin{problem}\label{pr:gauss}
By machine, \(c_4 = 0.143191337023821\ldots\), and the composite \(\Ac_{(4)}(1) \to \Ac_{(4)}(1) \to \Ac_{(4)}(1)\) of two lifts \(\Theta_{(4)}\) is multiplication by \(c_4^2/16 = 0.00128148493741\ldots\).
\end{problem}

In a model with several classes a form is a vector indexed by classes, a lift has a kernel matrix \(K\), the class weights form a diagonal matrix \(D\), the lift is \(v \mapsto KDv\), and \(\langle z, v \rangle := z^t D v\).

\begin{problem}\label{pr:transpose-toy}
For \(K_0 = \bigl(\begin{smallmatrix} 1&2\\ 3&-1 \end{smallmatrix}\bigr)\), \(K_1 = \bigl(\begin{smallmatrix} 2&1\\ 0&-1 \end{smallmatrix}\bigr)\), \(D_0 = \diag(1, \tfrac12)\), \(D_1 = \diag(\tfrac13, 1)\), \(D_2 = \diag(1, \tfrac14)\), \(v = (2,1)^t\), \(z = (1,-2)^t\): \(K_1D_1K_0D_0v = (\tfrac{15}2, -\tfrac{11}2)^t\), \(z^tD_2K_1D_1K_0D_0v = \tfrac{41}4\), and with both lifts transposed onto \(z\), \(K_0^tD_1K_1^tD_2z = (\tfrac{31}6, -\tfrac16)^t\), whose pairing \((\cdot)^tD_0v\) with \(v\) is again \(\tfrac{41}4\).
\end{problem}

Each datum \(\phi_j(\xi)\) is a pure tensor \(\bu \otimes \phi_f\) of a real Fock vector \(\bu\) and a finite datum \(\phi_f\) (the integrated Heisenberg action of the proof).  In the proof: the Weil actions, the join and \(\Heis(w)\) respect this separation, a translate staying pure with any scalar absorbed in its real vector; on the real factor the characters act through their weights, which do not depend on the tame character; and the finite join \(\boxtimes\) on tensors of finite data is injective, with a linear left inverse \(\boxtimes^{-1}\).

(ii) \emph{Archimedean zero term}, in the proof: for a translate \(\Heis(w)\) and a real Fock vector \(\bu\) there is a real Fock vector \(\bu_0\) with \(\zt(\Heis(w)(\bu \otimes \phi_f)) = \bu_0 \otimes \zt_f(\Heis_f(w_f)\phi_f)\) for every \(\phi_f\); here \(\zt_f\) is the finite zero term and \(w_f\) the finite coordinate of the translation in the Weil model (the finite part of \(w\), conjugated into the frames and divided by the finite idele of the Weil data).  For the mixed model see \cite{MVW}, \cite{We}.

\begin{definition}\label{def:evaluated}
For \(0 \le i \le \varrho\), with \(A\) and \(\zt\) those of pair \(i\), the \emph{evaluated pair datum} is
\begin{gather*}
  \phi^{\ev}_i(\xi, y) := \boxtimes^{-1}\,\zt_f\,\Heis_f(A_f)\bigl(\mathbf 1_{2i+1}(\xi)^y \boxtimes \mathbf 1_{2i+2}(\xi)^y\bigr), \\
  \text{where } \mathbf 1_j(\xi)^y := \omega_j(y_{j-1}, y_j)\,\mathbf 1_j(\xi) ,
\end{gather*}
with \(A_f\) as in (ii); it is a tensor of finite data on \(\Hom(X_{2i}, X_{2i+1})\) and on \(\Hom(X_{2i+1}, X_{2i+2})\).
\end{definition}

In the proof, \(\boxtimes\phi^{\ev}_i(\xi, y)\) is the finite zero term itself.  Put \(\phi_j(\xi)^{g,y} := \omega_j((g_{j-1}, y_{j-1}), (g_j, y_j))\,\phi_j(\xi)\).

\begin{problem}\label{pr:separate}
For the fixed pair data and harmonics there is a real Fock vector \(\bu_i\), depending on the outer points only through \(g_{2i}, g_{2i+1}, g_{2i+2}\), with \(\zt\,\Heis(A)(\phi_{2i+1}(\xi)^{g,y} \boxtimes \phi_{2i+2}(\xi)^{g,y}) = \bu_i \otimes {\boxtimes}\,\phi^{\ev}_i(\xi, y)\) for every tame \(\xi\) and all finite points \(y\).
\end{problem}

(iii) \emph{Definite theta lifts}, in the proof.  For a form \(\psi\) of the prescribed weight on \(X\) and a linear form \(l\) on its values put \(\coeff_l\psi(\gamma) := l(\gamma_\infty^{-1}\psi(\gamma_f))\), \(\gamma_\infty\) acting on the values.  For a family of definite lifts \(\Theta^X_o\), \(o \in O\), between fixed spaces of algebraic forms of the prescribed source and target weights, with Weil actions \(\omega_o\), the coefficient property is: for every polynomial \(P\) on \(\Hom(X, X')_\infty\), with Fock vector \(v_P\), and every \(l\) there is \(l'\) such that for all \(o\), finite data \(\phi_f\), forms \(\psi\) and \(h \in U(X')(\A)\),
\[
  \int_{[U(X)]} \theta\bigl(\omega_o(\gamma, h)(v_P \otimes \phi_f)\bigr)\,\coeff_l\psi(\gamma)\,d\gamma = \coeff_{l'}\bigl(\Theta^X_o(\phi_f)\psi\bigr)(h) .
\]
In the proof the lifts of the definite chain, \(O\) the tame characters, have it.  See \cite{Gr}, \cite{KV}, \cite{Ho}.

\begin{problem}\label{pr:closure}
Let \(\Lambda \colon \mathsf L \to \mathsf V\) be \(\mathbb C\)-linear from a finite-dimensional complex vector space to a Hausdorff complex topological vector space, and \(I \colon \mathsf H \to \mathsf V\) continuous on a topological space.  If \(I(u) \in \Lambda(\mathsf L)\) for \(u\) in a dense subset of \(\mathsf H\), then for all \(u \in \mathsf H\).
\end{problem}

For polynomials \(P\) on \(\Hom(X,X')_\infty\) and \(Q\) on \(\Hom(X',X'')_\infty\), let \(PQ\) be their product after transposing the variables of \(P\) into the \(X\)-block and placing those of \(Q\) in the \(X''\)-block.  In the proof, for finite data \(v, v'\) on the two edges, \((v_P\otimes v)\boxtimes(v_Q\otimes v') = v_{PQ}\otimes(v\boxtimes v')\), and on such joins the triple action at \((\gamma,\gamma',h)\) acts on the two factors by the rank-twisted Weil actions of \(X\to X'\) at \((\gamma,\gamma')\) and of \(X'\to X''\) at \((\gamma',h)\).

\begin{problem}\label{pr:poly}
Fix automorphic characters \(\mu_o\), \(\mu'_o\) of the norm-one torus, \(o \in O\), and families \(\Theta^X_{1,o} \colon X \to X'\), \(\Theta^X_{2,o} \colon X' \to X''\) with the same intermediate space of forms, having (iii) for the rank twists of \(\mu_o\) and \(\mu'_o\); let \(\omega^\triangle_o\) be their triple action.  For \(P\), \(Q\) as above and every \(l\) there is \(l''\) such that for all \(o\), finite data \(\phi_f\), \(\phi'_f\), forms \(\psi\) on \(X\) and \(h \in U(X'')(\A)\),
\begin{multline*}
  \int_{[U(X)]} \coeff_l\psi(\gamma) \int_{[U(X')]} \theta\bigl(\omega_o^\triangle(\gamma, \gamma', h)(v_{PQ} \otimes (\phi_f \boxtimes \phi'_f))\bigr)\,d\gamma'\,d\gamma \\
  = \coeff_{l''}\bigl(\Theta^X_{2,o}(\phi'_f)\,\Theta^X_{1,o}(\phi_f)\,\psi\bigr)(h) .
\end{multline*}
\end{problem}

In the proof the polynomial vectors are dense in the real Fock model in its \(L^1\) norm (for the Hilbert space see \cite{Ba} and Problem~\ref{fm-pr:monomials}); the families of double integrals, with pointwise topology in the other arguments, depend continuously on \(\bu\); and for fixed \(l\) the families \(\coeff_{l''}(\cdots)\) form a finite-dimensional space as \(l''\) varies.  Every polynomial is a sum of products \(PQ\), and both sides are linear in the finite tensor.  For these same two families the proof gives the following coefficient formula: for every real Fock vector \(\bu\) and every \(l\) there is \(l''\) with
\[
  \int_{[U(X)]} \coeff_l\psi(\gamma) \int_{[U(X')]} \theta\bigl(\omega_o^\triangle(\gamma, \gamma', h)(\bu \otimes {\boxtimes}\Phi)\bigr)\,d\gamma'\,d\gamma = \coeff_{l''}\bigl(\Theta^X_o(\Phi)\psi\bigr)(h)
\]
for all \(o\), finite pair tensors \(\Phi\), forms \(\psi\) and \(h\), where \(\Theta^X_o(v \otimes v') := \Theta^X_{2,o}(v')\Theta^X_{1,o}(v)\), extended linearly.  In the chain, \(l''\) is chosen before the tame character.  In the proof, for pair \(i\) the characters \(\mu\), \(\mu'\) of \eqref{eq:two-periods} are those of lifts \(2i+1\) and \(2i+2\), the families are \(\Theta^X_{1,\xi} = \Theta^X_{\xi,2i+1}\) and \(\Theta^X_{2,\xi} = \Theta^X_{\xi,2i+2}\) with the rank twists of these characters, \(O\) the tame characters, and \(\omega^\triangle\) of \eqref{eq:two-periods} is their triple action.


\section{Induction on pairs}\label{sec:induction}

For \(0 \le j < m\), a form \(\psi\) on \(X_j\) and a linear form \(l\), let \(\Wc_{l,\psi}\) on \(R_j(\A)\) be formed as \(\Wc_{\lambda,\varphi}\) in Section~\ref{ss-sec:theorem}: \(\Wc_{l,\psi}(n\gamma) = \psi_\Uc(n)\coeff_l\psi(\gamma)\) for \(n \in \Uc(\A)\), \(\gamma \in U(X_j)(\A)\), with \(\psi_\Uc(n) := \psi_E\bigl(\tfrac12 j_E\operatorname{tr}(\st_{j+1}^*\st_{j+1}n)\bigr)\) for every \(n \in \Uc(\A)\).

\begin{problem}\label{pr:sourcefn}
Let \(\len Y_j\) be odd and \(\Wc\) be equivariant on \(R_j(\A)\) with \(\Wc(\gamma) = \coeff_l\psi(\gamma)\) for all \(\gamma \in U(X_j)(\A)\).  Then \(\Wc = \Wc_{l,\psi}\).
\end{problem}

For finite pair tensors \(\Phi_0, \dots, \Phi_i\), \(\Theta^X_\xi(\Phi_i \otimes \dots \otimes \Phi_0)\) is the composite of the first \(i+1\) double lifts of the definite chain, pair \(j\) acting through \(v \otimes v' \mapsto \Theta^X_{\xi,2j+2}(v') \circ \Theta^X_{\xi,2j+1}(v)\).

\begin{theorem}[two lifts at a time]\label{thm:two-lifts}
Fix \(\lambda\), archimedean points \(g_k\) of \(U(Y_k)\) \((k \ge 0)\) and \(0 \le i \le \varrho\).  There is a linear form \(l\) on the values of forms on \(X_{2i+2}\), independent of the tame character, of \(\varphi\) and of the finite points, such that for every tame character \(\xi\), every \(\varphi \in S\), all finite points \(y_k \in U(Y_k)(\A_f)\), and every \(h\) in \(U(X_{2i+2})(\A)\) if \(i < \varrho\), in \(U(Y_m)(\A)\) if \(i = \varrho\),
\[
  \Wc_{2i+2}\bigl(h\,(g_{2i+2}, y_{2i+2})\bigr) = \coeff_l\Bigl(\Theta^X_\xi\bigl(\phi^{\ev}_i(\xi, y) \otimes \dots \otimes \phi^{\ev}_0(\xi, y)\bigr)\varphi\Bigr)(\bar h) .
\]
Here \(h\) acts on \(Y_{2i+2}\) through its definite part, and \(\bar h = h\) if \(i < \varrho\), \(\bar h = \diag(h, 1) \in U(X_m)(\A)\) if \(i = \varrho\).
\end{theorem}

\paragraph{What has been proved.}  On the definite group the translated stage function after each pair is a coefficient of the definite chain, with one linear form for all tame characters, source forms and finite points.  The problems prove: the isotropy of the stacked map (Problems~\ref{pr:isotropic} and~\ref{pr:commute}), the product of two constant terms (Problem~\ref{pr:ct-product}), the separation of the archimedean variables (Problem~\ref{pr:separate}), the coefficient formula for polynomial vectors with its extension by density (Problems~\ref{pr:poly} and~\ref{pr:closure}), and that an equivariant function on an odd string is determined by its values on the definite group (Problem~\ref{pr:sourcefn}).  The hints indicate how Theorem~\ref{thm:two-lifts} follows from these.  Taken as given: the product formula~(i), the archimedean zero term~(ii), the coefficient property of the definite lifts~(iii), the integration over a Bessel group~(iv), the two-period identity~\eqref{eq:two-periods} and the coefficient formula for all real Fock vectors (both stated from the proof; the problems give their ingredients), and the statements marked ``in the proof''.  Chapter~\ref{s3-chap} needs Definitions~\ref{def:stage} and~\ref{def:evaluated}, Problem~\ref{pr:stage} and Theorem~\ref{thm:two-lifts}: for ramified source forms it applies support bounds and quotient integration to periods whose dependence on the archimedean points and the finite level classes separates.

\begin{leanremark}
\raggedright Theorem~\ref{thm:two-lifts} specializes \tl{Data.exists\_unfoldStage\_eq\_algFormCoeff} \cite{GL}, which allows any double-lift data with the prescribed finite evaluation for pair \(i\) and an independent final archimedean point; here they are the middle or last pair and \(g_{2i+2}\).  The stage functions are \tl{unfoldStageOn} and \tl{unfoldStage}, the stacked map \tl{stackedStringMap}, the evaluated pair datum \tl{evaluatedDatum}.  Identity \eqref{eq:two-periods} specializes \tl{DoubleLiftData.exists\_integral\_constTermAtStringMap\_mul\_eq} (the action is \tl{tripleWeil}); the display before Section~\ref{sec:induction} is \tl{integral\_theta\_tripleWeil\_eq\_lift\_comp}.  Input (i) is \tl{theta\_arrowProduct}, (ii) specializes \tl{exists\_zeroTerm\_rho\_integrated\_eq}, (iii) is \tl{IsThetaLift}, proved for the chain by \tl{isThetaLift\_thetaKernel}, and (iv) is \tl{isCocompact\_ratBessel} with \tl{integral\_besselMeasure\_eq\_of\_mean}.  The other statements marked ``in the proof'' are: \tl{exists\_unfoldStageOn\_even\_eq\_sourceFn} and \tl{formAdjoint\_stackedStringMap\_mul\_stackedStringMap} (pairs, isotropy); \tl{isEquivariant\_unfoldStageOn} (equivariance of the stages); \tl{transposedLiftY\_apply}, \tl{integral\_mul\_liftY} and \tl{isRapid\_conjFormY} (the transpose); \tl{Heisenberg.Rep.mul} and \tl{homBicharacter\_eq} (the Heisenberg law); \tl{arrowProduct} and \tl{homHeisenbergOp\_arrowProduct} (the join); \tl{mul\_constTermAtStringMap\_unipotent\_mul} (the first edge); \tl{datumY}, \tl{pairWeil\_integrated\_all}, \tl{boxProd\_integrated}, \tl{CoordChange.map\_integrated}, \tl{exists\_rho\_integrated\_eq}, \tl{hasWeight\_liftChar}, \tl{injective\_arrowTensor} and \tl{leftInverse\_arrowTensor} (the separation); \tl{DoubleLiftData.arrowTensor\_finEval} (the evaluated datum); \tl{arrowProduct\_integrated\_polyVector} and \tl{arrowProduct\_twistedPairWeil} (polynomial joins); \tl{Fock.denseRange\_polyVector} (density), the remaining steps of the coefficient formula being in the proof of \tl{integral\_theta\_tripleWeil\_eq\_lift\_comp}; and \tl{exists\_unfoldStage\_eq\_algFormCoeff\_of\_sourceFn} (the order of \(l''\) and \(\xi\) in the chain).  Automorphy at stage \(j+1\) is part of \tl{ConjFormY}.  In Problem~\ref{pr:stage} absolute convergence is \tl{integrable\_unfoldStage\_mul} and the equality is \tl{integral\_unfoldStageOn\_mul\_transposedLiftY}.\par
\end{leanremark}


\section*{Hints}
\subsection*{\S\ref{sec:stage}}
\begin{description}
\item[\ref{pr:stage}] Move \((g_j, y_j)\) into the datum and apply Theorem~\ref{s1-thm:transfer} to \(\Wtr_j\).
\end{description}
\subsection*{\S\ref{sec:strings}}
\begin{description}
\item[\ref{pr:A-instance}] The adjoint of \(\binom{P}{Q}\) into \(V^- \oplus V''\) is \((-P^*, Q^*)\).
\item[\ref{pr:adjoint-degree}] The forms pair grades \(g\) and \(-g\).
\item[\ref{pr:isotropic}] Compute both products on each \(f'_q\), using \(\bar c_V = c_V\).
\item[\ref{pr:last-pair}] Problem~\ref{da-pr:piN}.
\end{description}
\subsection*{\S\ref{sec:two-ct}}
\begin{description}
\item[\ref{pr:ct-product}] As in Problem~\ref{pr:adjoint-degree}, \(a \mapsto a^*\) preserves positive degree; use (i) and Fubini on the compact groups.
\end{description}
\subsection*{\S\ref{sec:definite}}
\begin{description}
\item[\ref{pr:gauss}] Problem~\ref{tn-pr:U1U1}; keep track of the class weight in each lift; in Python, sum \(\alpha^4 e^{-\pi\lvert\alpha\rvert^2}\) over \(\lvert\Re\alpha\rvert, \lvert\Im\alpha\rvert \le 8\).
\item[\ref{pr:separate}] Split data and points into real and finite factors; apply (ii).
\item[\ref{pr:closure}] Finite-dimensional subspaces of a Hausdorff topological vector space are closed.
\item[\ref{pr:poly}] Apply (i) to the definite triple \(X, X', X''\), with its own scalar \(z^X_\theta\); after Fubini use (iii) over \(\gamma\), then over \(\gamma'\) for \(z^X_\theta Q\).
\end{description}
\subsection*{\S\ref{sec:induction}}
\begin{description}
\item[\ref{pr:sourcefn}] Problem~\ref{pr:odd}.  For Theorem~\ref{thm:two-lifts}: apply the coefficient formula with \(\bu = (c_Rc_{R'}/z_\theta)\bu_i\), and Problem~\ref{pr:sourcefn} at the intermediate even stages.
\end{description}

\begin{thebibliography}{MVW}
\bibitem[Ba]{Ba} V.~Bargmann, \emph{On a Hilbert space of analytic functions and an associated integral transform, Part~I}, Comm.\ Pure Appl.\ Math.\ \textbf{14} (1961), 187--214.
\bibitem[Bo]{Bo} A.~Borel, \emph{Introduction aux groupes arithm\'etiques}, Hermann, Paris, 1969.
\bibitem[GL]{GL} L.~Sernau, \emph{GRH for Dirichlet \(L\)-functions, in Lean~4}, a shorter rewriting of the Lean proof, text by Claude, \url{https://github.com/lsernau/grh-lean}, commit \texttt{72fead1}, readable edition \texttt{grh/}.
\bibitem[Gr]{Gr} B.~H.~Gross, \emph{Algebraic modular forms}, Israel J.\ Math.\ \textbf{113} (1999), 61--93.
\bibitem[Ho]{Ho} R.~Howe, \emph{Transcending classical invariant theory}, J.\ Amer.\ Math.\ Soc.\ \textbf{2} (1989), 535--552.
\bibitem[KV]{KV} M.~Kashiwara and M.~Vergne, \emph{On the Segal--Shale--Weil representations and harmonic polynomials}, Invent.\ Math.\ \textbf{44} (1978), 1--47.
\bibitem[MVW]{MVW} C.~M\oe glin, M.-F.~Vign\'eras and J.-L.~Waldspurger, \emph{Correspondances de Howe sur un corps \(p\)-adique}, Lecture Notes in Math.\ 1291, Springer, 1987.
\bibitem[We]{We} A.~Weil, \emph{Sur certains groupes d'op\'erateurs unitaires}, Acta Math.\ \textbf{111} (1964), 143--211.
\end{thebibliography}

\end{document}
